% CPSY 1291 -- Recitation handout 4
% Compile:  latexmk -lualatex handout-04-gradients-chain-rule.tex
\documentclass[11pt]{article}
\usepackage{cpsy1291-handout}

\begin{document}

\handouthead{4}{One neuron, and the calculus of learning}
  {Angles, gratings, derivatives, gradients, and the chain rule}
  {Week 4 \ $\cdot$\ hold this after Lecture 6; Assignment 2 is due Tue 10/13}

\section*{What this session is for}
\addcontentsline{toc}{section}{What this session is for}

Assignment 2 is about what a single neuron computes and how it learns. Two
things it assumes are not taught in lecture, and this is where you get them.

\textbf{Part A} is geometry: the dot product written as
$\|\boldsymbol{w}\|\|\boldsymbol{x}\|\cos\theta$, angles in radians, and how to
build a sinusoidal grating --- the stimulus used to measure orientation tuning in
real visual cortex and, in the assignment, in a network.

\textbf{Part B} is calculus: derivatives, partial derivatives, the gradient, the
chain rule, and gradient descent. Every learning rule in this course is one
line of calculus applied to an error. Once you can differentiate a squared
error, the delta rule stops being something to memorize and becomes something
you can rederive on the exam in thirty seconds.

If you have had a calculus course, Part B is revision and you can skim it. If
you have not, read it slowly --- it is written for you, and nothing beyond it is
required all term.

\begin{keybox}
\ask{The one idea.} A derivative answers exactly one question: \emph{if I nudge
this input a little, how much does the output move, and in which direction?}
Everything else --- gradients, backpropagation, the whole training loop --- is that
question asked about more variables at once.
\end{keybox}

\part*{Part A --- The geometry of one neuron}
\addcontentsline{toc}{section}{Part A --- The geometry of one neuron}

\section{The dot product, written two ways}

A linear neuron computes $a = \boldsymbol{w}\cdot\boldsymbol{x} + b$. The same
quantity can be written

\[
\boldsymbol{w}\cdot\boldsymbol{x} \;=\; \sum_i w_i x_i
\qquad\text{and}\qquad
\boldsymbol{w}\cdot\boldsymbol{x} \;=\; \|\boldsymbol{w}\|\,\|\boldsymbol{x}\|\,\cos\theta .
\]

The first tells you how to compute it. The second tells you what it \emph{means},
and it separates the response into two independent causes:

\begin{itemize}
  \item $\|\boldsymbol{x}\|$ --- how strong the input is (its contrast, its
        overall firing rate);
  \item $\cos\theta$ --- how well its \emph{pattern} matches the neuron's
        preferred pattern.
\end{itemize}

\begin{keybox}
\ask{Why this matters.} A unit can respond strongly either because the stimulus
matched its preference or because the stimulus was simply intense. Those are
different scientific claims. Normalizing every stimulus to unit length removes
the second cause, and then the response is $\|\boldsymbol{w}\|\cos\theta$ --- a
pure measure of pattern match.
\end{keybox}

\begin{lstlisting}
xn = x / np.linalg.norm(x)                 # unit length: intensity removed
cos = (w @ xn) / np.linalg.norm(w)         # in [-1, 1]
theta_deg = np.degrees(np.arccos(np.clip(cos, -1, 1)))
\end{lstlisting}

The \texttt{np.clip} is not paranoia. Floating point returns things like
$1.0000000000000002$, and \texttt{arccos} of that is \texttt{nan}.

\subsection{Which stimulus drives a unit hardest?}

Among all inputs of a fixed length, the response
$\|\boldsymbol{w}\|\|\boldsymbol{x}\|\cos\theta$ is largest when $\cos\theta=1$
--- that is, when $\boldsymbol{x}$ points in the same direction as
$\boldsymbol{w}$. \textbf{The optimal stimulus for a linear unit is its own
weight vector.} This is why plotting a unit's weights as an image is a
meaningful thing to do, and it is the entire basis for calling the first layer
of a network ``Gabor-like''.

\section{Radians, degrees, and gratings}

\subsection{Angles}

NumPy's trigonometric functions take \textbf{radians}. A full turn is $2\pi$
radians, $180^\circ = \pi$ radians.

\begin{lstlisting}
np.deg2rad(45)      # 0.7853...   use these two functions and never
np.rad2deg(np.pi)   # 180.0       convert by hand
\end{lstlisting}

Orientation, unlike direction, is defined modulo $180^\circ$: a bar tilted at
$10^\circ$ and one at $190^\circ$ are the same bar. When you average or compare
preferred orientations, work modulo $180$, or a unit preferring $175^\circ$ and
one preferring $5^\circ$ will look maximally different when they are
$10^\circ$ apart.

\subsection{A sinusoidal grating in four lines}

A grating is a stripe pattern: brightness varying as a sine wave along one
direction. It has an orientation $\theta$, a spatial frequency $f$ (stripes per
image), and a phase $\varphi$ (where the stripes sit).

\begin{lstlisting}
def grating(size=11, theta=0.0, freq=2.0, phase=0.0):
    y, x = np.mgrid[0:size, 0:size] / size          # coordinates in [0, 1)
    proj = x * np.cos(theta) + y * np.sin(theta)    # distance along theta
    return np.sin(2 * np.pi * freq * proj + phase)
\end{lstlisting}

The middle line is the idea: project each pixel's position onto the direction
$\theta$, then make brightness a sine of that projection. Everything
perpendicular to $\theta$ has the same projection, which is why the stripes come
out parallel.

\subsection{Tuning curves, and the number to report}

A \textbf{tuning curve} is a unit's response plotted against a stimulus
parameter. Measure it by sweeping the parameter and recording the response.

Phase is a nuisance parameter: a unit that likes vertical stripes may respond
strongly to one phase and weakly to its opposite, purely because the light and
dark bars swapped. The standard fix is to average over several phases at each
orientation, which is what makes the curve about orientation and not about
where the stripes happened to land.

The usual summary of how sharp a tuning curve is:

\begin{keybox}
\ask{Half-width at half-maximum (HWHM).} Find the peak. Find the two
orientations either side where the response has fallen to halfway between the
peak and the baseline. Half the distance between them is the bandwidth. Smaller
means more sharply tuned. Macaque V1 simple cells sit near $23^\circ$.
\end{keybox}

\part*{Part B --- The calculus of learning}
\addcontentsline{toc}{section}{Part B --- The calculus of learning}

\section{A derivative is a slope}

$f'(x)$ --- also written $\dfrac{df}{dx}$ --- is the slope of $f$ at $x$: how much
$f$ changes per unit change in $x$.

\begin{center}
\begin{tabular}{@{}lll@{}}
\toprule
$f(x)$ & $f'(x)$ & where you meet it \\
\midrule
$c$ (a constant)      & $0$                & biases that are not being trained \\
$x$                   & $1$                & everywhere \\
$x^2$                 & $2x$               & squared error \\
$e^{x}$               & $e^{x}$            & exponentials, softmax \\
$\log x$              & $1/x$              & cross-entropy, surprisal \\
$\sigma(x) = \frac{1}{1+e^{-x}}$ & $\sigma(x)\,(1-\sigma(x))$ & sigmoid units \\
$\tanh x$             & $1 - \tanh^2 x$    & RNNs \\
$\mathrm{ReLU}(x) = \max(0,x)$ & $1$ if $x>0$, else $0$ & almost every network \\
\bottomrule
\end{tabular}
\end{center}

Two remarks on that table. The sigmoid and $\tanh$ derivatives are written in
terms of the function's \emph{own output}, which is why a network can compute
them cheaply during backpropagation --- it already has the output. And ReLU has
no derivative at exactly $0$; every library simply picks $0$ and nothing
depends on the choice.

\begin{keybox}
\ask{The two combining rules, and they are all you need.}
Sum: $(f+g)' = f' + g'$.\quad
Chain: if $y = f(g(x))$ then $\dfrac{dy}{dx} = f'(g(x))\cdot g'(x)$.
\end{keybox}

\section{The chain rule as a pipeline}

Read the chain rule as bookkeeping about a pipeline rather than as an
identity. Suppose

\[
z = \boldsymbol{w}\cdot\boldsymbol{x} + b, \qquad
a = \mathrm{ReLU}(z), \qquad
L = \tfrac{1}{2}(a - y)^2 .
\]

To ask how the loss responds to $w_i$, walk backwards and multiply the local
slopes:

\[
\frac{\partial L}{\partial w_i}
= \underbrace{(a - y)}_{\partial L / \partial a}
\cdot \underbrace{\mathbb{1}[z>0]}_{\partial a / \partial z}
\cdot \underbrace{x_i}_{\partial z / \partial w_i}.
\]

Each factor is local: it involves only one stage and quantities that stage
already has. That locality \emph{is} backpropagation --- the algorithm is nothing
more than this multiplication, done once per layer, from the loss backwards.

\begin{keybox}
\ask{Where dead ReLUs come from.} If $z \le 0$ the middle factor is $0$, so the
whole product is $0$ and the weight does not move --- no matter how wrong the
answer was. A unit that outputs $0$ for every input in the training set can
never recover. You will see exactly this when you look inside the networks that
fail to learn XOR.
\end{keybox}

\section{Partial derivatives and the gradient}

With several inputs, a \textbf{partial derivative}
$\partial f/\partial w_i$ is the slope in the $w_i$ direction with every other
variable held fixed. Stack them and you have the \textbf{gradient}:

\[
\nabla_{\boldsymbol{w}} f =
\Big(\frac{\partial f}{\partial w_1}, \dots, \frac{\partial f}{\partial w_p}\Big).
\]

\begin{keybox}
\ask{The gradient points uphill.} It is the direction of steepest \emph{increase},
and its length says how steep. To decrease a loss you therefore step in
$-\nabla$. Every minus sign in every learning rule in this course is that
sentence.
\end{keybox}

Two gradients worth knowing by sight, because they turn a page of algebra into
one line:

\[
\nabla_{\boldsymbol{w}}\,(\boldsymbol{w}\cdot\boldsymbol{x}) = \boldsymbol{x},
\qquad
\nabla_{\boldsymbol{w}}\,\tfrac{1}{2}\|\boldsymbol{w}\|^2 = \boldsymbol{w}.
\]

\section{Gradient descent, and where the delta rule comes from}

The algorithm, in one line, repeated:

\[
\boldsymbol{w} \leftarrow \boldsymbol{w} - \eta\,\nabla_{\boldsymbol{w}} L .
\]

Now take a linear neuron $\hat{y} = \boldsymbol{w}\cdot\boldsymbol{x}$ and the
squared error $L = \tfrac{1}{2}(y - \hat{y})^2$, and turn the handle:

\[
\frac{\partial L}{\partial \boldsymbol{w}}
= (\hat{y} - y)\cdot \frac{\partial \hat{y}}{\partial \boldsymbol{w}}
= (\hat{y} - y)\,\boldsymbol{x}
\qquad\Longrightarrow\qquad
\Delta \boldsymbol{w} = \eta\,(y - \hat{y})\,\boldsymbol{x}.
\]

\begin{keybox}
\ask{That is the delta rule}, and also Rescorla--Wagner, and also the LMS rule
in signal processing. Three literatures, three names, one derivation ---
\emph{move the weights in proportion to the error, along the input}. If you can
reproduce these two lines you never have to memorize any of the three.
\end{keybox}

\subsection{The learning rate is the whole difficulty}

\begin{center}
\begin{tabular}{@{}ll@{}}
\toprule
$\eta$ too small & loss falls, slowly and smoothly; you run out of patience \\
$\eta$ about right & loss falls fast, then flattens \\
$\eta$ too large & loss oscillates, or jumps to \texttt{nan} within a few steps \\
\bottomrule
\end{tabular}
\end{center}

Plot the loss on a \textbf{log} $y$-axis (\texttt{plt.semilogy}). Progress that
looks like a flat line on a linear axis is often a clean straight descent on a
log axis, and the difference decides whether you conclude ``it is not learning''
or ``it needs more epochs''.

A decaying learning rate, $\eta_t = \eta_0/(1+t)$, is the classical choice: it
shrinks fast enough to settle but slowly enough to still travel any finite
distance. That is the Robbins--Monro condition, and it is worth remembering
when the fashionable schedulers show up later in the term.

\section{Checking a gradient you derived by hand}

The single most useful debugging tool in this part of the course. A derivative
is a limit of a difference quotient, so approximate it numerically and compare:

\begin{lstlisting}
def numerical_grad(f, w, eps=1e-5):
    g = np.zeros_like(w, dtype=float)
    for i in range(w.size):
        wp, wm = w.astype(float).copy(), w.astype(float).copy()
        wp[i] += eps
        wm[i] -= eps
        g[i] = (f(wp) - f(wm)) / (2 * eps)
    return g

# analytic vs numerical -- these should agree to ~1e-7
np.abs(my_grad(w) - numerical_grad(loss, w)).max()
\end{lstlisting}

Use the \emph{two-sided} difference shown here, not $(f(w+\epsilon)-f(w))/\epsilon$:
it is far more accurate for the same $\epsilon$. It is slow --- one pass per
parameter --- so run it once on a tiny example, confirm the match, then never
again.

\begin{keybox}
\ask{If the check fails.} A factor of exactly $2$ means a missing
$\tfrac{1}{2}$ in the loss. A sign flip means you wrote $(y - \hat{y})$ where
the derivative gives $(\hat{y} - y)$. Agreement everywhere except at a few
entries usually means a ReLU sitting exactly at its kink.
\end{keybox}

\section{Why gradient descent can fail}

For a linear neuron with squared error the loss surface is a bowl: one minimum,
and descent finds it from anywhere. The moment you add a hidden layer this stops
being true, and three things become possible.

\begin{itemize}
  \item \textbf{Local minima} --- a valley that is not the deepest valley.
        Descent stops there because every direction is uphill.
  \item \textbf{Saddle points} --- uphill in one direction, downhill in another,
        with a near-zero gradient in between. Progress stalls without the loss
        being low.
  \item \textbf{Dead units} --- the ReLU case above: a gradient of exactly zero,
        permanently.
\end{itemize}

\begin{keybox}
\ask{The practical consequence.} With a non-convex loss, \emph{the same network
trained twice with different random initializations can give different answers,
and one of them can simply fail}. This is not a bug to fix. It means a single
training run is a single observation, and a claim about an architecture needs
several seeds --- the same discipline that Recitation 3 applied to $t$-SNE.
\end{keybox}

\section{Exercises}

\ask{1.} Differentiate $L = \tfrac{1}{2}(y - wx)^2$ with respect to $w$, treating
$x$ and $y$ as constants. Then write the gradient-descent update.

\ask{2.} A unit computes $a = \mathrm{ReLU}(\boldsymbol{w}\cdot\boldsymbol{x})$
and currently has $\boldsymbol{w}\cdot\boldsymbol{x} = -3$ while the target is
$a = 5$. What is $\partial L/\partial \boldsymbol{w}$, and what happens on the
next update?

\ask{3.} You normalize every input patch to unit length and find that a unit's
response now correlates $0.98$ with $\cos\theta$ between the patch and the
weights, where before normalizing it correlated $0.31$. Explain, in one
sentence, using $\|\boldsymbol{w}\|\|\boldsymbol{x}\|\cos\theta$.

\ask{4.} Your analytic gradient is exactly twice your numerical gradient
everywhere. What is the bug?

\ask{5.} A training loss looks like a flat line for 500 epochs. Name two
different things this could be, and the one plot that distinguishes them.

\ask{6.} Two units have preferred orientations of $175^\circ$ and $5^\circ$. A
classmate reports their difference as $170^\circ$. What is wrong, and what is
the right answer?

\subsection*{Answers}

\ask{1.} $\dfrac{dL}{dw} = -(y - wx)\,x$, so
$w \leftarrow w + \eta\,(y - wx)\,x$. The sign flips because we step against the
gradient.

\ask{2.} The ReLU factor $\mathbb{1}[z>0]$ is $0$, so the whole gradient is
$\boldsymbol{0}$ and nothing changes --- no matter how large the error. The unit
is dead for this input.

\ask{3.} Before normalizing, the response confounded two causes,
$\|\boldsymbol{x}\|$ and $\cos\theta$, and variation in patch intensity
dominated. Fixing $\|\boldsymbol{x}\| = 1$ removes the first cause, leaving the
response proportional to $\cos\theta$ alone.

\ask{4.} A missing factor of $\tfrac{1}{2}$ in the loss: you differentiated
$(y-\hat{y})^2$ but your analytic expression is for
$\tfrac{1}{2}(y-\hat{y})^2$, or the reverse.

\ask{5.} (i) The learning rate is far too small and it is descending
imperceptibly; (ii) it is genuinely stuck --- dead units, or a saddle. Plot the
loss on a log axis: a steady straight descent means (i), a truly flat line means
(ii).

\ask{6.} Orientation is defined modulo $180^\circ$, so the difference is
$\min(170, 180-170) = 10^\circ$. The two units have nearly the same preference.

\section*{Cheat sheet}
\addcontentsline{toc}{section}{Cheat sheet}

\begin{center}
\begin{tabular}{@{}p{0.44\textwidth}p{0.5\textwidth}@{}}
\toprule
\textbf{You want} & \textbf{You write} \\
\midrule
length of a vector                 & \texttt{np.linalg.norm(x)} \\
unit-length version                & \texttt{x / np.linalg.norm(x)} \\
angle between two vectors          & \texttt{np.degrees(np.arccos(np.clip(c,-1,1)))} \\
degrees $\leftrightarrow$ radians  & \texttt{np.deg2rad} / \texttt{np.rad2deg} \\
orientation difference             & \texttt{d = abs(a-b) \% 180; min(d, 180-d)} \\
a grating                          & \texttt{np.sin(2*np.pi*f*proj + phase)} \\
gradient-descent step              & \texttt{w -= eta * grad} \\
delta rule                         & \texttt{w += eta * (y - yhat) * x} \\
loss curve worth reading           & \texttt{plt.semilogy(losses)} \\
check a hand-derived gradient      & two-sided difference, \texttt{eps=1e-5} \\
\bottomrule
\end{tabular}
\end{center}

\end{document}
